\documentclass[11pt]{article}

% ---------------------------------------------------------------------------
% Everlasting Options on HyperEVM --- project write-up / research note
% Build:  pdflatex everlasting-options-hyperevm.tex   (run twice for refs)
%    or:  tectonic everlasting-options-hyperevm.tex
% Self-contained: only stock TeX Live packages, no shell-escape needed.
% ---------------------------------------------------------------------------

\usepackage[margin=1in]{geometry}
\usepackage{amsmath,amssymb}
\usepackage{booktabs}
\usepackage{array}
\usepackage{enumitem}
\usepackage{listings}
\usepackage{xcolor}
\usepackage{microtype}
\usepackage{parskip}
\usepackage[colorlinks=true,linkcolor=black,citecolor=black,urlcolor=blue!55!black]{hyperref}

\lstdefinestyle{code}{%
  basicstyle=\ttfamily\footnotesize,
  breaklines=true,
  columns=fullflexible,
  keepspaces=true,
  showstringspaces=false,
  frame=single,
  rulecolor=\color{black!25},
  backgroundcolor=\color{black!4},
  xleftmargin=4pt,
  xrightmargin=4pt,
}
\lstset{style=code}

\newcommand{\code}[1]{\texttt{\small #1}}

\title{\textbf{Everlasting Options on HyperEVM}\\[2pt]
\large Funding-Settled, Liquidation-Free Options with On-Chain HyperCore Cover}
\author{SvennM\\[2pt]\small\url{https://github.com/your-handle/hyperevmlearning}}
\date{July 3, 2026}

\begin{document}
\maketitle

\begin{abstract}
\noindent
This note documents an in-progress, fully on-chain \emph{everlasting options} venue built on
Hyperliquid's HyperEVM (chain 998). Everlasting options never expire; in place of an expiry date,
longs pay an hourly funding fee equal to $(\text{mark}-\text{intrinsic})$ into a collateral pool,
turning the expiry cliff into a continuous carry (White \& Bankman-Fried, Paradigm 2021). The venue
is peer-to-pool, fully collateralized, and replaces liquidation with deterministic auto-settlement.
Slice~1 --- an everlasting put with a cumulative-funding engine and auto-settle --- is deployed and
smoke-tested on HyperEVM testnet. Slice~2 (a capped call spread, a protocol fee, and an opt-in yield
float) and Slice~3a (an \emph{uncapped} covered call with an on-chain cover ledger and cover-gate)
are code-complete behind a fuzz/invariant test suite that asserts solvency conservation on every
run. The paper's empirical core is a live-testnet \emph{spike} (Slice~3b) that drives a real HYPE
perpetual ``cover'' from a smart contract through Hyperliquid's CoreWriter system contract, pinning
several undocumented encoding rules and surfacing a blocking result: CoreWriter's action set exposes
no way to set leverage or margin mode, forcing a contract-opened perp into liquidatable isolated
margin --- which inverts the cover design from a perpetual hedge back to spot. We close with a
verified feasibility study for extending the same funding machine to real-world-asset (HIP-3) perps.
\end{abstract}

\tableofcontents
\vspace{1em}
\hrule

% ===========================================================================
\section{Introduction}
% ===========================================================================

\textbf{Everlasting options} are the option analogue of the perpetual swap. Rather than expire, a
position accrues a periodic \emph{funding} payment; longs pay shorts an amount that tracks the gap
between the option's traded mark and its current intrinsic value, so the contract can be held
indefinitely while remaining anchored to fair value~\cite{paradigm}. This removes the operational
pain of expiries and roll dates, and it turns an option book into something that looks, mechanically,
like a funding-settled perpetual.

\textbf{Why HyperEVM.} Hyperliquid pairs a high-performance on-chain central limit order book
(``HyperCore'', the L1) with a general EVM (``HyperEVM'') that can \emph{read} L1 state through
precompiles and \emph{write} L1 actions through a system contract (``CoreWriter''). That coupling is
unusual: an EVM smart contract can price an option against the same oracle Hyperliquid uses for its
own funding and liquidation, and can hedge or collateralize itself by placing real orders on the L1
order book --- without an off-chain relayer in the critical path. This project asks how far a
self-collateralizing, self-hedging options venue can be pushed purely on-chain on that substrate.

\textbf{Contributions.} This write-up consolidates the project's research to date:
\begin{enumerate}[leftmargin=1.4em,itemsep=2pt]
  \item A generalized \code{EverlastingMarket} contract (put and capped call) with a cumulative
        funding index, full collateralization, and auto-settlement instead of liquidation
        (Section~\ref{sec:arch}).
  \item An \emph{uncapped} covered-call design whose short-call risk is backed by an on-chain
        \emph{cover} position rather than cash, guarded by a cover-gate invariant
        (Section~\ref{sec:covered}).
  \item A live-testnet spike that operates a real HYPE perp cover from a contract via CoreWriter,
        pinning undocumented write/read encoding rules and identifying a leverage-control gap that
        forces a redesign (Section~\ref{sec:spike}).
  \item A verified feasibility study for everlasting options on HIP-3 real-world-asset perps
        (Section~\ref{sec:hip3}).
\end{enumerate}

\textbf{Scope and honesty.} This is a testnet research project, not audited or production software.
Collateral is a mock USDC; marks are posted by an off-chain keeper under on-chain guards; deployed
addresses are throwaway. Claims below are annotated as \emph{deployed}, \emph{code-complete}, or
\emph{empirically verified on testnet/mainnet} accordingly.

% ===========================================================================
\section{Background}
% ===========================================================================

\subsection{Everlasting option mechanics}

For a strike $K$ and underlying spot $S$, the intrinsic value per unit is
\[
\text{intrinsic}_{\text{put}} = \max(K-S,\,0), \qquad
\text{intrinsic}_{\text{call}} = \max(S-K,\,0).
\]
A keeper posts a \emph{mark} $m \geq \text{intrinsic}$ each funding period. Longs pay funding
\[
\text{funding per unit per period} \;=\; m - \text{intrinsic},
\]
accrued through a \emph{cumulative funding index} so that a position's owed funding is the difference
of the index between open and settlement times, scaled by size. Because a long's downside is bounded
(an option is never worth less than zero) and the pool is fully collateralized, the venue can replace
liquidation with \emph{auto-settlement}: when a position's accrued funding (and, for the covered call,
its mark losses) would exceed its posted collateral, anyone may permissionlessly settle it, capping
the loss at the collateral and returning the position to the pool.

\subsection{HyperEVM $\leftrightarrow$ HyperCore}

Two facilities connect the EVM to the L1:
\begin{itemize}[leftmargin=1.4em,itemsep=2pt]
  \item \textbf{Read precompiles} (address range \code{0x...0800}) expose L1 state to EVM
        \code{staticcall}: \code{oraclePx} / \code{markPx} (per-asset oracle and mark price),
        \code{position} (a user's perp position), \code{withdrawable}, \code{accountMarginSummary},
        and \code{coreUserExists}.
  \item \textbf{CoreWriter} (system contract \code{0x3333...3333}, entrypoint \code{sendRawAction})
        submits L1 actions: place limit order (action id 1), transfer between spot/perp, spot send,
        set account abstraction (unified account), add API wallet, and vault transfer --- 16 actions
        in total.
\end{itemize}
A recurring subtlety, pinned empirically in Section~\ref{sec:spike}, is that \emph{reads and writes
do not share a number system}: the price/size scale used by CoreWriter orders differs from the scale
returned by the read precompiles, and for builder-deployed (HIP-3) markets even the \emph{asset id}
differs between the read and write paths.

% ===========================================================================
\section{System architecture}\label{sec:arch}
% ===========================================================================

The contracts are Foundry/Solidity \code{0.8.35} (\code{evm\_version = cancun}) and depend on the
\code{hyper-evm-lib} submodule, which wraps the HyperCore precompiles and CoreWriter. The
implementation plan was hardened by a 45-agent adversarial audit (16 findings, 4 blockers fixed)
before any contract was written; the resulting fixes are cited inline in the code as
\code{AUDIT F1..F6}, and later review findings as \code{I1..I3}.

\subsection{Oracle --- \code{OracleLib}}
\code{OracleLib.spotWad()} reads HYPE's canonical price via \code{oraclePx(135)} and rescales it to
WAD (18-dp). The read precompile returns price $\times\,10^{6-\text{szDecimals}}$; for HYPE
($\text{szDecimals}=2$) that is $\times\,10^{4}$, so the library multiplies by $10^{14}$ to reach
WAD. This is the same feed Hyperliquid uses for its own funding and liquidation, consumed
trustlessly on-chain.

\subsection{Market --- \code{EverlastingMarket(side, K, W)}}
One contract serves both sides via \code{enum Side \{ PUT, CALL \}} and immutables $K$ (strike) and
$W$ (maximum payout per unit). Intrinsic value is clamped to $[0, W]$: a put uses $W = K$; a capped
call (a call \emph{spread}) uses $W = K_{hi}-K$, bounding payout at $W$.

\begin{itemize}[leftmargin=1.4em,itemsep=3pt]
  \item \code{openLong(qty)} escrows initial margin $\text{IM} = \text{qty}\cdot W$ of USDC from the
        pool's free balance into a locked balance, and guards $\text{mark}>0$, mark staleness
        (\code{MAX\_MARK\_AGE = 7200\,s}), and one position per trader.
  \item \code{postMark(m)} is keeper-only and enforces $m \geq \text{intrinsic}$, $m \leq W$, a
        deviation cap (\code{MAX\_MARK\_DEV\_BPS = 2000}, i.e.\ 20\%), and staleness. It advances
        \code{cumFunding += (m - lastIntrinsic) * periods} with \code{FUNDING\_PERIOD = 3600\,s}.
        Funding uses the \emph{contemporaneous} intrinsic (fix \code{F3}); a stale gap is treated as
        a recoverable re-seed rather than a jump (fix \code{F5}).
  \item \code{close()} settles funding plus mark PnL and releases escrow \emph{before} paying out, so
        the payout \code{require} can never false-revert (fix \code{F2}); a loss is floored at the
        trader's collateral (``never below zero'').
  \item \code{settle(t)} is the permissionless auto-settle path when funding alone would exceed
        collateral.
  \item \textbf{Protocol fee:} \code{protocolFeeBps} (capped at \code{MAX\_FEE\_BPS = 2000}) skims a
        cut of the funding carry \emph{after} the trader is paid, floored at the free pool so it can
        never cause a shortfall.
  \item \textbf{Yield float:} an opt-in \code{IYieldAdapter} sweeps idle free capital to yield while
        keeping at least \code{reserveBps = 2000} (20\%) liquid in-contract; \code{deployedToYield}
        tracks principal.
\end{itemize}

\subsection{Uncapped covered call --- \code{CoveredCallMarket}}\label{sec:covered}
The ``real product'' is an \emph{uncapped} everlasting call: $\text{intrinsic} = \max(S-K,0)$ with
\emph{no upper clamp}. Uncapped upside for the buyer cannot be cash-collateralized cheaply, so the
short-call risk is instead backed by a long-underlying \emph{cover}:
\begin{itemize}[leftmargin=1.4em,itemsep=3pt]
  \item \textbf{Cover ledger.} \code{coverQty} and a quantity-weighted average \code{coverEntry};
        \code{increaseCover(qty)} updates the entry as
        $\text{coverEntry} \leftarrow \dfrac{\text{coverQty}\cdot\text{coverEntry} + \text{qty}\cdot S}
        {\text{coverQty}+\text{qty}}$, and \code{reduceCover(qty)} reverts if it would drop cover
        below the net written size.
  \item \textbf{Cover-gate.} \code{openLong} requires
        $\text{coverQty} \geq \text{netWritten} + \text{addQty}$ --- a call can be written only if the
        pool already holds enough cover to back it.
  \item \textbf{Settlement predicate.} Because initial margin here is premium-only
        ($\text{qty}\cdot\text{mark}$, not $\text{qty}\cdot W$), \code{settle(t)} uses the \emph{full}
        net-loss predicate $\text{netLoss} = \text{markLoss} + \text{funding} - \text{markGain}$
        (fix \code{I2}); a funding-only predicate would miss mark losses and lock out the keeper.
\end{itemize}
The open question this design defers is \emph{realizing the cover into USDC on payout} --- i.e.\
whether the pool can actually hold a live cover position on HyperCore. That is exactly what the
Slice-3b spike (Section~\ref{sec:spike}) probes.

\subsection{Off-chain keeper}
A TypeScript keeper posts the mark hourly. The mark is a geometric basket of Black--Scholes prices
across twelve horizons, approximating the everlasting's continuously-rolled term structure:
\[
\text{mark}_{\text{put}} = \sum_{i=1}^{12} 2^{-i}\,\text{BS}_{\text{put}}(K,\tau_i),
\qquad
\text{mark}_{\text{call}} = \sum_{i=1}^{12} 2^{-i}\big[\text{BS}_{\text{call}}(K,\tau_i)-\text{BS}_{\text{call}}(K_{hi},\tau_i)\big],
\]
clamped to $[\text{intrinsic}, W]$ before posting. Volatility $\sigma$ is currently hand-set
(default 0.9). The keeper is the one trusted off-chain component; its authority is bounded on-chain by
the \code{postMark} guards (bounds, deviation cap, staleness).

% ===========================================================================
\section{Verification}
% ===========================================================================

Testing is Foundry-based and mixes unit, fuzz, invariant, and live-fork tests --- 46 tests at the
Slice-2 snapshot, growing to roughly 84 test functions across 24 files on the current branch with
Slice-3a/3b added. Three tests are live-fork (they read the testnet oracle and need an RPC env).

The invariant suites are the load-bearing ones. Each uses \emph{nonce prediction} so that the fuzz
handler \emph{is} the keeper (it can post marks the way production would), and each has an
\code{afterInvariant()} non-vacuity guard (\code{assertGt(marksPosted, 0)}) so a suite cannot pass by
doing nothing. Table~\ref{tab:inv} lists what each asserts.

\begin{table}[h]
\centering
\begin{tabular}{@{}lll@{}}
\toprule
\textbf{Suite} & \textbf{Invariant} & \textbf{Asserts} \\
\midrule
\code{EverlastingPut}  & \code{invariant\_solvency}     & $\text{bal} = \text{poolFree}+\text{poolLocked}+\sum\text{collateral}$ \\
\code{EverlastingCall} & \code{invariant\_solvency}     & same, with spot fuzzed above $K_{hi}$ \\
\code{FeeFloat}        & \code{invariant\_conservation} & $\text{bal} = \text{poolFree}+\text{poolLocked}+\sum\text{coll}+\text{fee}$ \\
                       &                                & \emph{and} $\text{adapter.balance} \geq \text{deployedToYield}$ \\
\code{CoveredCall}     & \code{invariant\_conservation} & $\text{bal} = \text{poolUsdc}+\sum\text{collateral}$ \\
\code{CoveredCall}     & \code{invariant\_coverGate}    & $\text{coverQty} \geq \text{netWritten}$ \\
\bottomrule
\end{tabular}
\caption{Invariants. The put suite was verified non-vacuous over 256 runs / 128k calls
($\approx$15k \code{postMark}s + 21k opens/settles).}
\label{tab:inv}
\end{table}

\textbf{An honesty note worth recording.} The covered-call solvency invariant was \emph{rewritten
during review}: the original \code{invariant\_solvencyGivenCover} was found vacuous and replaced
(commit ``\emph{replace vacuous solvency invariant with conservation + cover-gate}'', fix \code{I1}).
The replacement tests explicitly document their own limit --- the cover-gate proves enough cover
\emph{exists}, but not that the pool can fund \code{close()} payouts by realizing that cover into
USDC, which is precisely the Slice-3b question. Naming what a test does \emph{not} prove is, in
practice, more valuable than the green check.

% ===========================================================================
\section{The CoreWriter cover spike (Slice 3b)}\label{sec:spike}
% ===========================================================================

\textbf{Question.} Can a HyperEVM contract hold a real HYPE perpetual as cover --- open, read, and
close it --- entirely on-chain through CoreWriter? Before writing the production
\code{CoveredCallMarket3b}, a deliberately \emph{throwaway} spike
(\code{src/spike/PerpCoverSpike.sol}, an \code{onlyOwner} probe) was deployed to testnet
(\code{0x8Bc3f8...D4A8}) and driven through a full round-trip.

\subsection{What worked}
\begin{itemize}[leftmargin=1.4em,itemsep=2pt]
  \item \textbf{Contract-driven perp lifecycle.} Opened 0.5 HYPE long at \$79.22 and closed it, all
        from contract \code{onlyOwner} calls routed through CoreWriter; fills confirmed via
        \code{userFills}, position read back via the \code{position} precompile.
  \item \textbf{Unified account from the contract.} \code{setAbstraction(self, 2)} flipped the
        contract's Core account to a unified account, so spot USDC directly margins the perp --- no
        explicit spot$\to$perp class transfer needed.
  \item \textbf{PnL settles to spot USDC.} On close, the balance went $120.000 \to 119.140$ USDC in
        the unified spot balance --- so realizing cover into USDC on a covered-call payout (the
        deferred \code{I3} step) is mechanically clean: close the perp, USDC lands in spot, then
        \code{spotSend} onward.
  \item \textbf{Rescue path.} \code{spotSendUsdc} returned the full 119.14 USDC to the funding user.
\end{itemize}

\subsection{Encoding, pinned empirically}
The genuinely hard-won, undocumented bits (Table~\ref{tab:scale}, Table~\ref{tab:assetid}):

\begin{table}[h]
\centering
\begin{tabular}{@{}lll@{}}
\toprule
\textbf{Quantity} & \textbf{Encoding} & \textbf{Example} \\
\midrule
Order \code{limitPx} (write) & human $\times\,10^{8}$ & \$85 $\to$ \code{85e8} \\
Order \code{sz} (write)      & human $\times\,10^{8}$ & 0.5 HYPE $\to$ \code{5e7} \\
\code{oraclePx}/\code{markPx} (read) & price $\times\,10^{6-\text{szDec}}$ & HYPE ($\text{szDec}=2$): $\times\,10^{4}$ \\
\code{position.szi} (read)   & coins $\times\,10^{\text{szDec}}$ & --- \\
\bottomrule
\end{tabular}
\caption{Write-scale $\neq$ read-scale. A first attempt that reused the \emph{read} price scale
($\text{price}\times 10^{4}$) as an order \code{limitPx} \textbf{silently 0-filled} --- CoreWriter
expired the order with no revert and no fill.}
\label{tab:scale}
\end{table}

\begin{table}[h]
\centering
\begin{tabular}{@{}lll@{}}
\toprule
\textbf{Market} & \textbf{Read precompile id} & \textbf{CoreWriter action asset id} \\
\midrule
Main perp dex (dex 0), e.g.\ HYPE & \code{135} (= perp index) & \code{135} (same) \\
HIP-3 builder dex, e.g.\ \code{xyz:GOLD} & \code{dex$\cdot$10000+idx} (10003) & \code{100000+dex$\cdot$10000+idx} (110003) \\
\bottomrule
\end{tabular}
\caption{On the main dex the read and write ids coincide; on a builder (HIP-3) dex they differ.
Passing the write id to a read precompile reverts (\code{PrecompileError}).}
\label{tab:assetid}
\end{table}

CoreWriter is also \emph{fire-and-forget and delayed}: a fill lands a Core block later, so reads
issued immediately after a write still show pre-fill state (``poll one block after the write'').

\subsection{Live round-trip}
\begin{table}[h]
\centering
\begin{tabular}{@{}lll@{}}
\toprule
\textbf{Step} & \textbf{Action} & \textbf{Result} \\
\midrule
Fund  & Core spot send 120 USDC & account funded \\
Open  & 0.5 HYPE long @ \$79.22, marketable IOC & filled, isolated $10\times$ \\
Read  & \code{position} precompile & szi 0.5, entry \$79.22, liqPx \$71.74 \\
Close & \code{reduceOnly} IOC @ \$75.68 (bid) & filled, PnL $\to$ spot USDC \\
Sweep & \code{spotSend} to user & 119.14 USDC returned \\
\bottomrule
\end{tabular}
\caption{Full contract-driven round-trip on testnet. Round-trip cost \$0.86
($120.00 \to 119.14$), dominated by the thin \emph{testnet} book spread (bought \$79.22 ask, sold
\$75.68 bid) --- not representative of mainnet.}
\label{tab:roundtrip}
\end{table}

\subsection{The blocking finding}
A contract-opened perp cover opens \textbf{isolated, at the asset's max leverage ($10\times$ on HYPE),
and is liquidatable} --- in the live run the liquidation price sat at \$71.74, about $-9.4\%$ from
entry, backed by only $\approx\$4.86$ of margin. Even with the account set to unified, the
\emph{position} opened isolated $10\times$.

Critically, \textbf{CoreWriter's action set (ids 1..16) contains no \code{updateLeverage}, no
cross/isolated toggle, and no add-isolated-margin action.} From a pure on-chain flow the cover
therefore \emph{cannot} be made cross, $1\times$, or otherwise non-liquidatable. This breaks the
covered-call premise directly: the cover must be fully-funded / non-liquidatable, or a HYPE drawdown
liquidates the cover and leaves the pool short a naked uncapped call --- insolvency.

\subsection{Design fork}
\begin{itemize}[leftmargin=1.4em,itemsep=3pt]
  \item \textbf{Option A --- spot HYPE cover (recommended).} Hold spot HYPE (CoreWriter spot buy,
        asset $10000+\text{spotIndex}$). No leverage, no liquidation, ever --- the textbook
        covered-call structure of owning the underlying. Ties up full notional and earns no funding,
        but is trustless and bulletproof; close $=$ sell spot $\to$ USDC.
  \item \textbf{Option B --- perp cover + off-chain leverage set.} Keep the perp (capital-efficient,
        earns funding), but since CoreWriter cannot set $1\times$/cross, use action~9 (add API
        wallet): the contract approves an agent EOA, and an off-chain keeper signs a one-time
        \code{updateLeverage(cross, 1$\times$)} L1 action for the contract account. Adds a trusted
        off-chain signer --- less trustless.
  \item \textbf{Option C --- perp cover + active margin management.} Keep isolated but have a keeper
        continuously keep the position far from its liquidation price. Highest ongoing
        trust/complexity; a bad-block or oracle gap can still liquidate. Weakest.
\end{itemize}
\textbf{Recommendation.} The spike \emph{inverts} the earlier ``perp cover'' decision. Given the
covered call needs a non-liquidatable cover and CoreWriter cannot set leverage, Option~A (spot
cover) is now the clean path; Option~B is viable only if funding income and capital efficiency are
worth an off-chain agent. Learning this from an \$0.86 probe rather than a production position is the
whole point of spiking first.

% ===========================================================================
\section{Extension: everlasting options on HIP-3 RWA perps}\label{sec:hip3}
% ===========================================================================

A separate feasibility study (verdict: \textbf{green --- both legs empirically verified}) asks whether
the same funding machine can be pointed at \emph{real-world-asset} perps: equities, indices,
commodities, and FX listed on HIP-3 builder markets. The dominant HIP-3 RWA venue is Trade.xyz
(\code{perp\_dex\_index = 1}, $\approx$98 markets; S\&P-500 licensed as of March 2026). Trading is
permissionless on-chain --- the $\sim$500k HYPE stake gates \emph{deploying} a market, not trading.

\textbf{Read path --- verified on mainnet.} A HyperEVM contract reads any \code{xyz:} market's oracle
via the \code{oraclePx}/\code{markPx} precompiles, with read index $=\text{dex}\cdot 10000+\text{idx}$.
On-chain reads matched the live info-API oracle across six markets (Table~\ref{tab:rwa}), confirming
the pricer and funding engine can consume the HIP-3 oracle on-chain.

\begin{table}[h]
\centering
\begin{tabular}{@{}lrrr@{}}
\toprule
\textbf{Market} & \textbf{Read id} & \textbf{Precompile \code{oraclePx}} & \textbf{Info-API oracle} \\
\midrule
\code{xyz:GOLD}  & 10003 & \$4116.6 & \$4116.2 \\
\code{xyz:NVDA}  & 10002 & \$192.85 & \$192.77 \\
\code{xyz:TSLA}  & 10001 & \$393.69 & \$393.66 \\
\code{xyz:AAPL}  & 10009 & \$307.45 & \$307.47 \\
\code{xyz:SP500} & 10052 & \$7443.7 & \$7442.1 \\
\code{xyz:ORCL}  & 10011 & \$139.60 & \$139.62 \\
\bottomrule
\end{tabular}
\caption{On-chain \code{staticcall} vs live info-API oracle (mainnet, 2026-07-02).}
\label{tab:rwa}
\end{table}

\textbf{Write/hedge path --- verified on testnet.} A deployed HyperEVM contract funded a HyperCore
account and sent a marketable IOC via CoreWriter (action asset id $=100000+\text{dex}\cdot 10000+\text{idx}$)
that \emph{filled}: position \code{szi = 0.5 felix:TEST1}, \code{entryPx = 43.31}. So a contract can
open, close, and hence delta-hedge a HIP-3 market on-chain.

\textbf{Risks to design around.} (i) A deployer-run centralized oracle (a $\sim$3\,s relayer,
clamped $\leq 1\%$ per update) lags fast moves and directly hits pricing and hedge accuracy; (ii) a
deployer \code{haltTrading} force-settles all positions to mark --- for an open-ended everlasting
option this needs a settlement fallback (secondary oracle / TWAP / auto-unwind); (iii) fees
$\approx 2\times$ native perps, isolated-margin only, and fire-and-forget delayed writes hit
slippage and rebalance latency; (iv) $>$90\% single-builder concentration is a systemic dependency.
Precedent exists (Synapse Labs' Hypercall runs options on a HIP-3-only SpaceX underlying), but the
general pattern is greenfield.

% ===========================================================================
\section{Roadmap}
% ===========================================================================

\begin{table}[h]
\centering
\renewcommand{\arraystretch}{1.2}
\begin{tabular}{@{}clp{6.3cm}@{}}
\toprule
\textbf{Slice} & \textbf{Adds} & \textbf{Status} \\
\midrule
1  & Everlasting put, isolated pool, funding, auto-settle & \textbf{Deployed} \& smoked on testnet (\code{0xD48a...C135}, $K=\$47$) \\
2  & Capped call spread, protocol fee, yield float & Code + tests complete; two-market testnet deploy pending \\
3a & Uncapped covered call: option accounting, cover ledger, cover-gate & Code + tests complete \\
3b & CoreWriter perp-cover spike & \textbf{Complete} --- findings invert the cover design (perp $\to$ spot) \\
4  & Partial margin + liquidation engine & Planned \\
5  & Insurance fund + ADL & Planned \\
--- & Multi-strike ladder; frontend; mainnet + external audit & Audit-gated \\
\bottomrule
\end{tabular}
\caption{Build slices. Only the Slice-1 \code{EverlastingPut} is live on-chain.}
\label{tab:roadmap}
\end{table}

% ===========================================================================
\section{Discussion}
% ===========================================================================

Three lessons generalize beyond this project. First, \textbf{spike before you build}: the most
consequential result here --- that on-chain leverage control does not exist in CoreWriter --- cost
\$0.86 and a weekend to find, and it changed a core architectural decision that would have been
expensive to unwind after the production contract was written. Second, \textbf{the interesting risk
lives in the encoding}, not the algebra: write-scale $\neq$ read-scale and read-id $\neq$ write-id are
undocumented, produce \emph{silent} failures (a 0-fill with no revert), and are exactly the kind of
thing only a live round-trip surfaces. Third, \textbf{invariant tests are only as good as their
non-vacuity}: replacing a vacuous solvency invariant with a conservation identity plus an explicit
cover-gate --- and documenting what the gate does \emph{not} prove --- was worth more than the
original green check.

The venue's remaining trust surfaces are named and bounded: an off-chain keeper posts marks (bounded
by on-chain guards), and any HIP-3 extension inherits the builder's oracle and halt authority (which
needs a settlement fallback). The near-term path is to build the Slice-3b covered call on a
\emph{spot} cover, finish the two-market Slice-2 deploy, and only then revisit the capital-efficiency
case for a perp cover with an off-chain leverage signer.

% ===========================================================================
\section{Conclusion}
% ===========================================================================

An everlasting-options venue can be assembled almost entirely on-chain on HyperEVM: it can price
against Hyperliquid's own oracle through read precompiles, collateralize and hedge itself by placing
real L1 orders through CoreWriter, and settle deterministically without a liquidation engine.
Slice~1 is live on testnet; slices 2 and 3a are code-complete under a solvency-conserving fuzz/invariant
suite; and the Slice-3b spike both proved contract-driven perp cover works and found the one missing
primitive --- on-chain leverage control --- that reshapes the cover design. The end goal is a
permissionless everlasting-options venue whose cover is real and fully on-chain, extensible to
real-world-asset perps whose read and write paths are already empirically verified.

% ===========================================================================
\appendix
\section{Deployed artifacts (testnet, chain 998)}
% ===========================================================================
\begin{itemize}[leftmargin=1.4em,itemsep=2pt]
  \item \code{EverlastingPut} --- \code{0xD48aabBd8ad161A68F471045E4c12D67c879C135}, $K=\$47$
        (Slice-1, live).
  \item \code{PerpCoverSpike} --- \code{0x8Bc3f867d40576485D8031D1a599A11a0a2FD4A8} (Slice-3b,
        throwaway probe).
  \item Toolchain: Foundry (solc \code{0.8.35}, \code{cancun}); submodules \code{forge-std},
        \code{hyper-evm-lib}; TypeScript keeper (\code{ethers}); GitHub Actions CI
        (\code{forge fmt --check}, \code{build --sizes}, \code{test}).
\end{itemize}

\section{Reusable CoreWriter facts}
\begin{itemize}[leftmargin=1.4em,itemsep=2pt]
  \item Order \code{px}/\code{sz} $=$ human $\times\,10^{8}$; marketable $=$ IOC through the book;
        poll reads one Core block after a write.
  \item A fresh contract Core account defaults to classic (disabled) + isolated max leverage;
        \code{setAbstraction(self, 2)} $\to$ unified.
  \item PnL on close realizes to the account's spot USDC balance.
  \item Main-dex order asset id $=$ read perp index (HYPE $=135$); the $100000+$ offset is
        builder-dex-only.
\end{itemize}

\section{Plain-language summary}
\emph{An options market that never expires, built on Hyperliquid. Instead of an expiry date, holders
pay a small hourly ``funding'' fee, so a position can be held forever while staying priced to fair
value. It is fully backed by collateral and settles automatically instead of liquidating. The put is
live on testnet; the call machinery is written and tested; and a weekend experiment proved a smart
contract can open and close a real HYPE position on Hyperliquid's order book on its own --- while
revealing that the on-chain toolkit has no way to set leverage, which sent the collateral design back
to simply owning the asset.}

% ===========================================================================
\begin{thebibliography}{9}
% ===========================================================================
\bibitem{paradigm} D.~White and S.~Bankman-Fried.
  \emph{Everlasting Options.} Paradigm, 2021.
  \url{https://www.paradigm.xyz/writing/everlasting-options}.
\bibitem{hldocs} Hyperliquid.
  \emph{HyperEVM $\leftrightarrow$ HyperCore precompiles, CoreWriter, Asset IDs, and HIP-3.}
  Hyperliquid documentation.
\bibitem{hyperevmlib} \emph{hyper-evm-lib} --- Solidity wrappers for HyperCore precompiles and
  CoreWriter. \url{https://github.com/hyperliquid-dev/hyper-evm-lib}.
\bibitem{tradexyz} Trade.xyz --- HIP-3 real-world-asset perpetuals on Hyperliquid.
\end{thebibliography}

\end{document}
